\section{Ordinary-angle C1 with arbitrary positive shape openings}
\label{sec:c1-arbitrary-openings}

The unit assumptions on $a,t,M,J$ in
Theorem~\ref{thm:c1-ordinary-angle} can all be removed. The large
positive packets must be retained before taking a bounded Gram residue.

\begin{theorem}[Ordinary angles and arbitrary positive openings]
Use the unchanged native C1 source (OA2), with $c=2$ and
$\lambda=n=0$. Assume only that $k,s=g/f$ are positive bounded units,
$\eta=f^2>0$ is infinitesimal, and
\[
\begin{gathered}
 \zeta=g^2=s^2\eta,\quad K=1+k^2,\quad
 R=K-ak^2\eta,\quad T=K-k^2s^2\eta t,\\
 a,t>0,\quad M=R-t>0,\quad J=T-a>0.
\end{gathered}\tag{AO1}
\]
Then the full original parent is tensor SOS at every independent
projective test and every real or complex state. No lower bound on
$a,t,M,J$, and no boundedness assumption on the original event
coefficients, is required.
\end{theorem}

Since $a<T<K$ and $t<R<K$, both shapes are bounded, and $R,T$ have
residue $K$. In particular
\[
       p=t/R\in(0,1),\qquad h=a/T\in(0,1)              \tag{AO2}
\]
are bounded at every opening rate. The seven native coquotients have
the same residues as in (OA3). All original events are positive:
$A/a$ and $\alpha/a$ are $\eta^2$ times positive units, while
$d/a$ is $\eta$ times a positive unit; positivity of the remaining
events also follows from the explicit bounds below.

\paragraph{Retain the four large packets exactly.}
Put $\Delta=s/k$, $\Lambda=RTc_3/(2a)$, and
\[
 H_2=4Tc_2,\quad H_b=4Tc_b,\quad H_3=4Rc_3,\quad H_5=4Rc_5.
\]
All four have residue $8K^2$. For a bounded finite testing angle define
\[
 D_u=1+u^2,\quad D_v=1+v^2,\quad
 x=X+uZ,\quad y=uW-Y,\quad
 q_A=v(W+uY),\quad q_B=v(Z-uX),
\]
\[
                    N=x-pq_A,\qquad L=y-hq_B.
\]
The four nonnegative packets are
\[
 \Pi_1=H_2\frac tM N^2,\quad
 \Pi_2=H_b\frac MtD_vN^2,\quad
 \Pi_3=\Delta^2H_3\frac JaD_vL^2,\quad
 \Pi_4=\Delta^2H_5\frac aJL^2.                         \tag{AO3}
\]
Their possibly unbounded pure source coefficients stay in their
separate square roots. They never enter the bounded remainder Gram.
The exact whole-parent identity is
\begin{align*}
 G:=\Delta^2\Lambda F-\sum_{i=1}^4\Pi_i
 ={}&H_2\{N^2-2q_AN+(1-p)q_A^2\}
     +H_bD_v\{N^2+2q_AN+pq_A^2\}\\
 &+\Delta^2H_3D_v\{L^2+2q_BL+hq_B^2\}
     +\Delta^2H_5\{L^2-2q_BL+(1-h)q_B^2\}
     +\mathcal W,                                      \tag{AO4}\\
 \mathcal W={}&4RT\Delta^2\eta^2c_4D_uD_vY^2
   +4RT\Delta^4k^6\eta^2c_\alpha D_u(vY-Z)^2\\
 &+2RT\Delta^2\eta c_dD_u\{(vY-Z)^2-D_vZ^2\}.
\end{align*}
This is the original four-packet identity (SD3)--(SD4), valid without
unit shapes. For example, the first subtraction uses only
$R=M+t$ and $x=N+(t/R)q_A$; its result is
$H_2\{N^2-2q_AN+(M/R)q_A^2\}$. The other three follow in the
same way from $R=M+t$ and $T=J+a$.

\paragraph{A bounded coordinate change with no mixed inverse.}
Let $w=(W+uY)/D_u$, which lies entirely in the testing field. Then
\[
\begin{gathered}
 q_A=D_uv w,\quad
 X=N+pD_uv w-uZ,\quad
 q_B=v(D_uZ-uN-puD_uv w),\\
 W=w+\frac{u(L+hq_B)}{D_u},\qquad
 Y=uw-\frac{L+hq_B}{D_u}.                             \tag{AO5}
\end{gathered}
\]
These are exact inverse formulas. Their only testing denominator is
the positive pure testing unit $D_u$; their source coefficients use
only $R^{-1},T^{-1}$. In particular no determinant involving $p,h,u,v$
is inverted. Set $\widehat G=D_uG$. It is a polynomial in
$N,L,u,w,v,Z$ with bounded coefficients, since this multiplication
clears the denominators in (AO5).

Suppose $v$ is infinitesimal and $u$ is a finite positive unit.
The exact polynomial $\widehat G$ has a bounded Gram on
\[
                    N,\ L,\ vw,\ vZ,\ \eta w,\ \eta Z.
                                                               \tag{AO6}
\]
After expanding only the squares in (AO4), every term contains a
product of two listed roots. The coefficient errors carry $v$, $\eta$,
or a native coquotient difference of zero residue. To display the
Gram residue, use the purely formal substitutions
$N=\eta\xi$, $L=\eta\chi$, $v=\eta\nu$ and read the coefficient of
$\eta^2$. It is
\[
 \mathcal G=16K^2D_u\xi^2+16\Delta^2K^2D_u\chi^2
                         +8\Delta^2K^3D_u^2S,        \tag{AO7}
\]
where
\[
 S=u^2w^2+s^2k^4Z^2+
   \frac{D_u}{K}\nu^2\left(\frac{k^2}{s^2}w^2+Z^2\right)
   -2Ku\nu wZ.                                       \tag{AO8}
\]
There is no shape parameter in this residue other than the fixed
units $k,s$. In particular no assertion that $v/\eta$ is bounded
has been made. Formula (AO7) follows directly from
$H_2,H_b,H_3,H_5\sim8K^2$, $c_4,c_\alpha\sim2K$, $c_d\sim4K^2$:
the $p$ and $h$ terms in the two scalar pairs sum to one, while the
two linear normal coefficients have zero residue.

\paragraph{Finite unit Alice angles: the same scalar, fewer hard roots.}
The capacity identity (OA10), with $W$ replaced by $w$, proves a
strict six-root Gram unless $\bar u=\bar k$. At the matched angle,
\[
 S=(\nu Z-kw)^2+
       \left(\frac{k}{s}\nu w-sk^2Z\right)^2.
\]
Define
\[
 r_1=vZ-\eta kw,\qquad r_2=(k/s)vw-\eta sk^2Z.
\]
The exact relation
\[
 wr_1-(s/k)Zr_2=\eta k(s^2Z^2-w^2)                    \tag{AO9}
\]
closes every state with $\overline{s^2Z^2-w^2}\ne0$ by adding its
identically zero squared relation to the Gram on $(N,L,r_1,r_2,\eta)$.
The multiplier in this relation is never inverted.

Normalize the testing state. If $N$ or $L$ has nonzero residue,
$\widehat G$ has positive residue and is already certified. Otherwise
the inverse formulas force $w$ or $Z$ to be a unit. When (AO9) has
zero multiplier residue, $Z$ is a unit; normalize $Z=1$.
The negative-core branch has $w>0$ by the pure testing core Schur
calculation printed in the ordinary-angle proof. Thus only
$\bar w=\bar s$ remains. The nonnegative-core branch and literal
axes are direct original-event SOS exits.

\paragraph{The exact contact with arbitrary shape rates.}
The native contact (OA4), under (AO5), is
\[
 (u_c,w_c,v_c,N_c,L_c)
       =\left(k,\frac{sR}{K},ks\eta,0,0\right).        \tag{AO10}
\]
All four packets (AO3) and all their first derivatives vanish there.
The original native contact identities consequently give exact value
and first-derivative zero for $\widehat G$ as well. This remains true
when a packet coefficient is unbounded: each coefficient is an element
of the source field multiplying a square that vanishes exactly.

Use the five roots
\[
       N,\ L,\ v-v_c,\ \eta(w-w_c),\ \eta(u-k).       \tag{AO11}
\]
The complete polynomial $\widehat G$ belongs to
$(N,L,v)^2+\eta(N,L,v)+(\eta^2)$, by (AO4)--(AO5).
Translation by $v_c=\eta ks$ preserves this inclusion. Its constant
normal coefficient is $\eta^2$ times a bounded polynomial in the
two soft deviations; its linear normal coefficients are $\eta$ times
such polynomials. Exact contact and stationarity make their constant
and the required first coefficients vanish. Finite monic Taylor
division therefore gives an exact bounded Gram on (AO11), just as in
(OA14), with no division by any testing/source comparison.
Higher Taylor terms have zero residue after this allocation.

The residue Gram is now particularly simple. Its first two blocks are
$16K^3$ and $16\Delta^2K^3$. Its remaining block is
$8\Delta^2K^5$ times the half-Hessian of $S$ at
$(\nu,w,u)=(ks,s,k)$ and $Z=1$. The latter is exactly (OA15), with
positive leading principal minors
\[
                         K,\qquad4k^4,\qquad12s^2k^6/K.
\]
Thus the entire five-root Gram is strict independently of all four
shape-opening rates.

\paragraph{The outer charts without bounded event weights.}
These charts admit a simpler uniform retained-packet comparison.
The exact source formulas give
\begin{align*}
 m/a&\ge\frac{2k^2c_2}{s^2Rc_3},&
 c/a&\ge\frac2T,\\
 D/a&\ge\frac{2c_5}{Tc_3}-\frac{\eta c_d}{c_3},&
 \beta/a&\ge\frac{2k^2c_b}{s^2Rc_3}-\frac{\eta c_d}{c_3}.
                                                               \tag{AO12}
\end{align*}
All four right sides are positive bounded units. Choose a positive
real $\mu$ smaller than their residues. In $F/a$, retain the four
nonnegative event packets with weights
$m/a-\mu,c/a-\mu,D/a-\mu,\beta/a-\mu$. The remaining complete
parent has the four corresponding weights all equal to $\mu$,
and the unchanged weights
\[
 d/a=\eta\,d_1,\quad A/a=\eta^2A_2,\quad
 \alpha/a=\eta^2\alpha_2,
\]
where $d_1,A_2,\alpha_2$ are positive units. This is an exact
decomposition at the present test, not a claim that the capped source
is globally physical.

If the homogeneous Bob sine has nonzero residue, the four event
amplitudes have trivial common state kernel, as proved in the
ordinary-angle theorem. Their bounded positive Gram absorbs the
infinitesimal signed core. If Bob is small and $u$ is infinitesimal,
use the six roots
$(x,y,vW,vZ,\eta uW,\eta Z)$ of the same theorem. The four hard
weights $\mu$ give a strict four-root block and the two weak diagonal
coefficients are positive units. The signed soft cross term has an
additional factor $u$, so its Gram coefficient has zero residue.
At the reciprocal Alice endpoint put $q=1/u$ and use
\[
                x'=qX+Z,\ y'=W-qY,\ vY,\ vX,
                       \eta Y,\ \eta qX.
\]
The hard residue is
$4\mu\{(vY-x')^2+x'^2+y'^2+(y'+vX)^2\}$;
the weak diagonals have positive unit coefficients. Assign the signed
soft cross term as a multiple of $q(vX)(\eta Y)$, which has zero
residue. These are the exact finite and reciprocal monomial
allocations already proved in the ordinary-angle theorem; they use
no shape inverse. Literal infinity follows in the homogeneous chart.

These outer charts and the finite-unit Alice analysis exhaust all
independent tests. Restore the retained packets, then the positive
pure factors $D_u$, $\Delta^2\Lambda$, or $a$ used in the respective
charts. Every radical and inverse is confined to one field. Complex
states split into their real and imaginary parts.

\paragraph{Remaining scope.}
This closes every positive opening $a,t,M,J$ on the ordinary-unit-$k$,
ordinary-unit-$s=g/f$, infinitesimal-$\eta$ row. It does not address
nonunit $k$ or $s$, non-small $\eta$, or literal zero openings that
are not defined by the native formulas. Actual-source entry retains
exactly the inherited contact-sign and nonzero-minor domain of (CE1).
